% Source-preserving integrated proof body; input from toe_round7_proofs.tex.
\subsubsection*{Round 7: the actual DET parent, its exact factorized gap, and the local-defect obstruction}
\label{sec:toe-r7-mirror}
\noindent\emph{The displayed arguments below are the proof-bearing source.  The current
repository regression is \veri{v1034\_\allowbreak parent\_\allowbreak mirror.py}; its finite checks audit the
implementation and do not replace the universal steps in the proof.}

\begingroup
\def\cH{\mathcal H}
\def\1{\mathbf 1}
\def\spec{\operatorname{spec}}
\def\ket#1{|#1\rangle}
\def\bra#1{\langle#1|}
\begin{keybox}[Scope and exact claim boundary]
The source audit separates three operators which had previously been discussed
under one ``mirror parent'' label.  The v1002 DET16 onsite term is the flat
defect projector $Q=\1-P_\phi$.  The v1008 wall scaffold instead tests the
four-flavour momentum-cell Hamiltonian
\[
 h(d)=d(N-2)+2Q,\qquad |d|\leq 3/5.
\]
The later v1027 identity constructs an additive dressed-number parent $N_a$;
it is not the same operator as $Q$.  For the actual v1008 momentum-factorized
Hamiltonian we prove, without a grid scan, the volume-independent many-body and
canonical-field threshold
\[
 \Delta_{\rm fact}\geq 1+\frac{\sqrt3}{2}=1.8660254037\ldots.
\]
This is the exact infimum of the continuous dispersion envelope; the source's
129-point grid instead attains $1.8660698881694362$.  It is not a real-space
locality theorem.  Indeed, for two real-space DET cells, a newly chosen bare
flavour-diagonal bond analogue $T_{\rm bare}$ satisfies (for even $n\geq4$)
\[
 \|T_{\rm bare}\ket{\Omega\Omega}\|^2=\frac n2|s|^2>0,
 \qquad (Q_x+Q_y)\ket{\Omega\Omega}=0.
\]
Consequently no finite $\kappa$ can obey
$T_{\rm bare}\geq-\kappa(Q_x+Q_y)$.  An explicitly gauge-dressed operator,
with the link shift applied as part of every hop, gives the same obstruction
inside finite $\mathbb Z_2/\mathbb Z_4$ Gauss sectors.  Thus the
desired elementary relative-$Q$ proof is impossible for the unmodified local
parent; a genuine small-perturbation theorem or a changed local bond is needed.
An explicit coupling-fixed bond repair is given, but is clearly marked as a
new, unlanded interaction.  No T3, T4, TOE, or RH claim is made.
\end{keybox}

\paragraph{Source contract and nomenclature.}

The certificate reads, but does not import, the following landed files at the
audited repository tip:
\begin{center}
\begin{tabular}{lll}
\toprule
source & exact object used here & status relevant to this note\\
\midrule
v1002 & $Q=\1-P_\phi$ on 16 complex modes & onsite projector only\\
v1008 & $h(d)=d(N-n/2)+JQ$, $n=4,J=2$ & momentum-cell scaffold\\
v1024 & positive seam/number-square parent & separate construction\\
v1027 & $N_a=\sum_i a_i^\dagger a_i$, $Q\leq N_a\leq nQ$ & algebraic identity\\
\bottomrule
\end{tabular}
\end{center}
In particular, ``DET32'' is sometimes used physically for 32 real components,
whereas the repository's operator is named DET16 because it has 16 complex
fermion modes.  We use $n$ for the number of complex modes and retain the
repository's DET16 name.

Let
\begin{equation}
 \ket{\Omega_\phi}=\frac{\ket0+e^{i\phi}\ket F}{\sqrt2},\qquad
 P_\phi=\ket{\Omega_\phi}\bra{\Omega_\phi},\qquad Q=\1-P_\phi.
 \label{eq:toe-r7-mirror-defect}
\end{equation}
Here $\ket F=c_1^\dagger\cdots c_n^\dagger\ket0$ and $n$ is even.  Both
$\ket0$ and $\ket F$ are assumed to be gauge singlets; this is the inherited
full-DET-singlet premise, not a conclusion of this note.  The common rotation
$e^{-i\phi N/n}$ removes $\phi$ from~\eqref{eq:toe-r7-mirror-defect} and commutes with $N$ and
ordinary flavour-diagonal hopping, so every spectrum below is phase independent.

The elementary algebra already confirms the proposed diagnosis:
\begin{equation}
 [Q,(-1)^N]=0,
 \qquad [Q,N]\neq0,
 \qquad \|[Q,N]\|_{\rm F}^2=\frac{n^2}{2}.
 \label{eq:toe-r7-mirror-paritynotnumber}
\end{equation}
Thus the flat projector preserves parity, not naive mirror number.  For the
v1008 analogue $n=4$, the last value is $8$; for DET16 it is $128$.

\paragraph{Exact theorem for the parent that v1008 actually computes.}

The v1008 routine \texttt{edge\_cluster\_hamiltonian} is
\begin{equation}
 h(d)=d\left(N-\frac n2\right)+JQ,\qquad
 n=4,\qquad J=2,\qquad |d|\leq t_{\rm src}=\frac35.
 \label{eq:toe-r7-mirror-hd}
\end{equation}
The scan uses $d=\pm t_{\rm src}\sin k$.  This is a separate Fock cell at every
sampled momentum, not a real-space hopping Hamiltonian.  In particular, a
standard symmetric nearest-neighbour term
$s\sum_x(c_x^\dagger c_{x+1}+\mathrm{h.c.})$ has dispersion magnitude
$2|s|$ (up to a momentum/phase convention).  Matching the v1008 dispersion
amplitude would therefore give $|s|=t_{\rm src}/2=3/10$, not $3/5$.

\paragraph{Lemma 1 (complete one-cell spectrum).}
For $1\leq r\leq n-1$, every occupation-$r$ state has energy
\begin{equation}
 E_r(d)=J+d\left(r-\frac n2\right).
 \label{eq:toe-r7-mirror-intermediate}
\end{equation}
On $\operatorname{span}\{\ket0,\ket F\}$, after removing $\phi$, the matrix is
\begin{equation}
 \begin{pmatrix}
 -nd/2+J/2&-J/2\\
 -J/2&nd/2+J/2
 \end{pmatrix},
\end{equation}
with eigenvalues
\begin{equation}
 E_\pm(d)=\frac J2\pm
 \sqrt{\frac{J^2}{4}+\frac{n^2d^2}{4}}.
 \label{eq:toe-r7-mirror-doublet}
\end{equation}

\paragraph{Proof.}
$P_\phi$ vanishes on occupation sectors $1,\dots,n-1$, which gives
\eqref{eq:toe-r7-mirror-intermediate}.  Its restriction to the empty/full doublet is the
rank-one matrix specified by~\eqref{eq:toe-r7-mirror-defect}; the displayed characteristic
polynomial gives~\eqref{eq:toe-r7-mirror-doublet}. \hfill$\square$

For the source values, $E_-(d)\leq0$, while the least intermediate energy is
$2-3/5=7/5>0$.  The ground state is therefore unique throughout the whole
declared dispersion interval.  A canonical $c_i$ or $c_i^\dagger$ maps this
ground doublet only to the $r=1$ or $r=n-1$ sectors, with nonzero amplitude.
The exact canonical spectral threshold at $x=|d|$ is consequently
\begin{equation}
 g(x)=\frac J2-\left(\frac n2-1\right)x
      +\sqrt{\frac{J^2}{4}+\frac{n^2x^2}{4}}.
 \label{eq:toe-r7-mirror-gx}
\end{equation}

\paragraph{Theorem 1 (source-fixed all-volume factorized gap).}
Let $K$ be any finite set of momentum cells, with arbitrary dispersions
$|d_k|\leq3/5$, and define the graded tensor-product Hamiltonian
$H_K=\sum_{k\in K}h(d_k)$.  Then $H_K$ has a unique product ground state and
\begin{equation}
 \operatorname{gap}(H_K)\geq
 \inf_{0\leq x\leq3/5}g(x)
 =1+\frac{\sqrt3}{2}.
 \label{eq:toe-r7-mirror-factgap}
\end{equation}
The same number is the uniform lower edge of canonical one-fermion spectral
support.

\paragraph{Proof.}
For $n=4,J=2$, differentiation of~\eqref{eq:toe-r7-mirror-gx} gives the unique interior
minimum
\begin{equation}
 x_*=\frac{J(n-2)}{2n\sqrt{n-1}}=\frac1{2\sqrt3}<\frac35,
\end{equation}
and substitution gives $g(x_*)=1+\sqrt3/2$.  The other empty/full eigenstate
costs at least $J=2$, and the $r=2$ sector also costs at least $2$, so neither
undercuts~\eqref{eq:toe-r7-mirror-factgap}.  The even cell Hamiltonians commute on the graded
tensor product and energies add.  Any state orthogonal to the product ground
excites at least one cell, proving the same bound independently of $|K|$.
\hfill$\square$

The analytic infimum $1.8660254037844386\ldots$ explains the lower envelope of
v1008's reported probe minimum; it is not the exact minimum of the finite scan.
The source grid $k_j=-\pi+j\pi/64$, $j=0,\ldots,128$, attains
$1.8660698881694362$.  The locality limit is equally exact: a projector local
in momentum is not a bounded-range real-space interaction.
Equation~\eqref{eq:toe-r7-mirror-factgap} therefore does not prove the missing local,
operator-valued TFPT parent theorem.

\paragraph{Why the desired relative-defect bound fails in real space.}

Put two DET cells on a bond and let
\begin{equation}
 D=J(Q_x+Q_y),\qquad
 T_{\rm bare}=s\sum_{i=1}^n
 (c_{x i}^\dagger c_{y i}+c_{y i}^\dagger c_{x i}).
 \label{eq:toe-r7-mirror-realbond}
\end{equation}
This is a bare-bond theorem.  The numerical choice $s=3/5$ below is a newly
chosen stress analogue, not a real-space coefficient derived from v1008.
For even $n\geq4$, the full-Fock CAR calculation contains only $2n$ nonzero
output monomials and does not require construction of the $2^{2n}$ matrix.

\paragraph{Lemma 2 (minimal kernel incompatibility).}
For every even $n\geq4$ and $s\neq0$,
\begin{align}
 D\ket{\Omega\Omega}&=0,
 &\bra{\Omega\Omega}T_{\rm bare}\ket{\Omega\Omega}&=0,\label{eq:toe-r7-mirror-zeroexp}\\
 \|T_{\rm bare}\ket{\Omega\Omega}\|^2&=\frac n2|s|^2>0,
 &(Q_x+Q_y)T_{\rm bare}\ket{\Omega\Omega}&=
 2T_{\rm bare}\ket{\Omega\Omega}.
 \label{eq:toe-r7-mirror-nonzeroimage}
\end{align}

\paragraph{Proof.}
Only the $\ket{F,0}$ and $\ket{0,F}$ components of the product DET vacuum can
hop.  Each has amplitude $1/2$.  The two hopping orientations and $n$ flavours
produce $2n$ mutually orthogonal states, each of modulus $|s|/2$, which proves
the norm in~\eqref{eq:toe-r7-mirror-nonzeroimage}.  Both sites have occupation $1$ or $n-1$
after the hop, hence both local projectors $P_x,P_y$ annihilate the image.
\hfill$\square$

The restriction $n\geq4$ is necessary for that counting statement.  At
$n=2$ the two orientations land on the same two basis states and add
constructively.  Direct CAR evaluation gives
$\|T_{\rm bare}\ket{\Omega\Omega}\|^2=2|s|^2$, rather than the former
formula's $|s|^2$.  The image is still nonzero, so the no-go below survives,
but the claimed universal norm formula did not.

\paragraph{Theorem 2 (no finite relative-defect constant).}
There is no finite $\kappa$ such that
\begin{equation}
 T_{\rm bare}\geq-\kappa(Q_x+Q_y).
 \label{eq:toe-r7-mirror-nokappa}
\end{equation}

\paragraph{Proof.}
If~\eqref{eq:toe-r7-mirror-nokappa} held, the positive operator
$B=T_{\rm bare}+\kappa(Q_x+Q_y)$ would have zero expectation in
$\ket{\Omega\Omega}$ by~\eqref{eq:toe-r7-mirror-zeroexp}.  Positivity then implies
$B\ket{\Omega\Omega}=0$.  Since $(Q_x+Q_y)\ket{\Omega\Omega}=0$, this would say
$T_{\rm bare}\ket{\Omega\Omega}=0$, contradicting~\eqref{eq:toe-r7-mirror-nonzeroimage}.
\hfill$\square$

For the newly chosen bare-bond analogue $n=4,J=2,s=3/5$,
\begin{equation}
 \|T_{\rm bare}\ket{\Omega\Omega}\|^2=\frac{18}{25}.
\end{equation}
Let $b=\sqrt{18/25}$ and normalize
$\ket w=T_{\rm bare}\ket{\Omega\Omega}/b$.  Direct CAR evaluation gives
$\bra wT_{\rm bare}\ket w=0$ and $\bra wD\ket w=2J$.  The compression of
$D+T_{\rm bare}$ to
$\operatorname{span}\{\ket{\Omega\Omega},\ket w\}$ is
\begin{equation}
 \begin{pmatrix}0&b\\b&2J\end{pmatrix}.
\end{equation}
Thus the variational principle gives the exact strict bound
\begin{equation}
 E_0(D+T_{\rm bare})\leq
 2-\frac{\sqrt{118}}5=-0.172556\ldots<0.
 \label{eq:toe-r7-mirror-ritz}
\end{equation}
The old product vacuum is not frustration free once ordinary hopping is on.
This does not prove gap closure; it proves that the proposed elementary
relative-$Q$ route cannot establish a gap.

\paragraph{Gauge-dressed theorem and actual link action.}
The bare theorem above makes no Gauss-sector claim.  For the $n=4$ analogue
and $q=2,4$, now adjoin an oriented link basis
$\{\ket e:e\in\mathbb Z_q\}$ and define
\begin{equation}
 U\ket e=\ket{e+1\bmod q},\qquad
 T_U=s\sum_{i=1}^n
 \left(c_{x i}^\dagger Uc_{y i}
 +c_{y i}^\dagger U^\dagger c_{x i}\right).
 \label{eq:toe-r7-mirror-gaugehopping}
\end{equation}
This operator is implementable on matter Fock space tensored with the finite
link space; it is fermion-even because each summand is a matter bilinear.
Take the endpoint Gauss constraints to be
$n_x-e=0$ and $n_y+e=0\pmod q$.  The starting vector
$\ket{\Omega\Omega}\ket{0}_E$ is physical because its site occupations are
$0$ or $4$.  Applying~\eqref{eq:toe-r7-mirror-gaugehopping}, including the link shift rather
than assigning a flux afterwards, produces only
\[
 (n_x,n_y,e)=(3,1,3\bmod q)\quad\hbox{or}\quad(1,3,1\bmod q),
\]
and in both cases $n_x-e=0$ and $n_y+e=0\pmod q$.  Hence the nonzero hopping
image is a physical flux excitation.  The checker applies the CAR and link
actions term by term for both $q=2$ and $q=4$ and verifies every nonzero output,
so Gauss preservation is not a post-hoc flux assignment.

Let $E^2\ket{\pm1}=\ket{\pm1}$ and
\begin{equation}
 H_0=J(Q_x+Q_y)+h_EE^2,\qquad h_E\geq0.
\end{equation}
Both terms annihilate the physical starting vector, whereas $T_U$ does not.
The same positivity-kernel argument therefore rules out any finite relative
bound of $T_U$ by this complete positive local parent.  Moreover, with
$\ket w=T_U\ket{\Omega\Omega,0_E}/b$, the exact Ritz compression is
\begin{equation}
 \begin{pmatrix}0&b\\b&2J+h_E\end{pmatrix},\qquad
 E_0(H_0+T_U)\leq \frac{2J+h_E}{2}
 -\sqrt{\frac{(2J+h_E)^2}{4}+b^2}.
 \label{eq:toe-r7-mirror-gaugeritz}
\end{equation}
Thus the electric cost is part of the excited diagonal entry.  For the same
newly chosen bond analogue $J=2,s=3/5$ and the illustrative choice $h_E=1$,
$b^2=18/25$ and~\eqref{eq:toe-r7-mirror-gaugeritz} is
$(25-\sqrt{697})/10=-0.140075756489\ldots<0$.

\paragraph{An exact local repair boundary.}

The following repair is included to make the result actionable, not to present
a new source claim.  On a bond let
\begin{equation}
 R_{xy}=\1-P_xP_y=Q_x+Q_y-Q_xQ_y.
\end{equation}
It obeys $0\leq R_{xy}\leq Q_x+Q_y$.  For the flavour-diagonal hopping
$K_{xy}$ with unit coefficient, $\|K_{xy}\|=n$ exactly.  Therefore
\begin{equation}
 R_{xy}K_{xy}R_{xy}\geq-nR_{xy}
 \geq-n(Q_x+Q_y).
 \label{eq:toe-r7-mirror-projectedbound}
\end{equation}
On a graph of maximum degree $z$, replacing ordinary hopping by its projected
version yields
\begin{equation}
 J\sum_xQ_x+s\sum_{\langle xy\rangle}R_{xy}K_{xy}R_{xy}
 \geq (J-zn|s|)\sum_xQ_x.
 \label{eq:toe-r7-mirror-rawrepair}
\end{equation}
For the newly chosen periodic-1D bond analogue
$J=2,z=2,n=4,s=3/5$, the margin is $-14/5$: projection alone is not enough.
This number is not source-derived.  Mode-normalizing $K/n$ would give the
positive margin $2-2(3/5)=4/5$, but that is another changed normalization.

There is also a coupling-fixed positive repair with no new tunable parameter:
\begin{equation}
 B_{xy}=|s|\left[nR_{xy}+\operatorname{sgn}(s)
 R_{xy}K_{xy}R_{xy}\right]\geq0.
 \label{eq:toe-r7-mirror-counterterm}
\end{equation}
For the newly chosen $s=3/5$ bond analogue its added coefficient is
$|s|n=12/5$ per bond; neither value is inferred from v1008.
The Hamiltonian $J\sum_xQ_x+\sum B_{xy}$ has the common product DET vacuum and
is bounded below by $J\sum_xQ_x$, hence has gap at least $J$.  If $P_x$ and
$K_{xy}$ are gauge invariant/covariant as assumed, so are $R_{xy}$ and
$B_{xy}$.  Equation~\eqref{eq:toe-r7-mirror-counterterm} is nevertheless a new many-body
bond counterterm and may alter the desired continuum dynamics.  It is an
explicit design option, not evidence that TFPT already contains it.

Keeping the ordinary hopping instead requires a real local-stability theorem.
The existing repository audit correctly notes that such theorems yield a
volume-independent window only under verified tensor-product/local-gap and
small-coupling premises, and that their unevaluated constants do not certify
the newly chosen $s=0.6$ bond point.  The exact obstruction above explains why
those premises cannot be replaced by the wished-for one-line inequality.

\paragraph{Reproducibility and claim boundary.}

The repository certificate \veri{v1034\_\allowbreak parent\_\allowbreak mirror.py} is read-only.  Its JSON separates 29 exact
algebra checks, six exact source-contract checks, and one floating-grid
regression.  The printed tag identifies each class.  The checker
performs an AST source-contract audit, exact $16\times16$ onsite algebra, exact
symbolic minimization, and a sparse CAR-monomial two-cell calculation.  It does
not build the $256\times256$ two-cell matrix; for general $n$ the obstruction
uses only $2n$ output monomials when even $n\geq4$.  It also retains the
exceptional $n=2$ collision as a regression and applies the $\mathbb Z_2$ and
$\mathbb Z_4$ link shifts explicitly.  The reported decimal radical values are
labelled evaluations; the separately labelled 129-point value is a floating
reproduction of the actual finite-grid minimum, not an exact trigonometric
certificate.

What is finished is precise: v1008's factorized parent has an exact all-volume
gap, while ordinary real-space hopping admits no relative bound against the
flat local DET defects.  What remains open is equally precise: no local
operator-valued TFPT Hamiltonian at a source-derived real-space coupling has
acquired a proved volume-uniform mirror gap, and no net-chiral T3/T4
construction follows.
\textbf{NO RH CLAIM.}
\endgroup
